\section{Migdal and nonadiabatic quantities}\label{app:E.18}

The drag mass of silicon has two values that must not be confused. The trace value, $14\,m_e$ or $\delta M/M = 2.73\times10^{-4}$ for natural silicon ($M = \SI{26.16}{GeV}$, \cref{sec:5.3}), is the sum of the drag-mass tensor over all nuclei and counts every electron; it is exact, but a kick does not measure it. A kick measures the on-site element, which the tables of \cref{sec:12.5} put at $28.7\,m_e$, $\delta M_{II}/M = 5.60\times10^{-4}$: $10.1\,m_e$ from the cores and $18.6\,m_e$ from the valence, with the band 12 to $24\,m_e$ for the valence and the sign of the off-site coupling not fixed. The other nonadiabatic inputs of \cref{sec:12.3,sec:12.5,sec:12.6} are likewise estimates from tables or models. The crystal Migdal rate is written there as a convolution of an electronic and a lattice response, and no single code has yet evaluated both pieces from first principles. A scalar that couples to nuclei out of proportion to their mass needs electronic matrix elements at reciprocal-lattice momenta. All the calculations below are at the leading nonadiabatic order of the Born--Oppenheimer parameter on the truncation label $\mathfrak{T}$ (\cref{sec:8.5}), harmonic ($n_A = 0$) and at zero temperature ($n_T = 0$). They differ in the many-body approximation $\mathcal{X}$, the lattice regime $\mathcal{K}$ and the Migdal regime $\mathfrak{M}$. \Cref{tab:E18-params} collects their parameters.

\subsection{Drag masses and inverse moments}\label{app:E.18-drag}

The electronic spectral measure of the force on nucleus $\kappa$ fixes every quantity of \cref{sec:12.5}. With $|0\rangle$ the clamped-nuclei ground state, $|a\rangle$ its excited states and $\Delta_a = E_a - E_0$, we define
\begin{equation}
  \dd\sigma_\kappa(\Delta) = \sum_a \bigl|\hat{\vect{q}}\cdot\langle a|\nabla_{\vect{R}_\kappa}V|0\rangle\bigr|^2\,\delta(\Delta - \Delta_a) ,
  \qquad
  \mathcal{I}_k \equiv \int \Delta^{-k}\,\dd\sigma_\kappa(\Delta) ,
  \label{eq:E.18-moments}
\end{equation}
where $V$ is the electron--nucleus potential. \Cref{sec:12.5} shows that the drag mass along $\hat{\vect{q}}$ is $\delta M_\kappa = 2\mathcal{I}_3$ and that the Migdal probability is $P_{\mathrm{Mig}} = q^2\mathcal{I}_4/M_\kappa^2$. Because $\Delta \ge \Eg$ on the support of $\dd\sigma_\kappa$, $\mathcal{I}_4 \le \mathcal{I}_3/\Eg$, and so $P_{\mathrm{Mig}} \le E_R\,\delta M_\kappa/(M_\kappa\Eg)$ with $E_R = q^2/2M_\kappa$. With the on-site drag mass $28.7\,m_e$ and the gap of \SI{1.17}{eV} this bound keeps $P_{\mathrm{Mig}}$ below $10^{-3}$ only up to $E_R \approx \SI{2.1}{eV}$ in silicon, and to less in germanium; the trace value $14\,m_e$ would give \SI{4.3}{eV}, but it does not bound the probability (\cref{sec:12.5-bounds}). Above these recoil energies the computed $\mathcal{I}_4$ is needed.

The three moments follow from two Sternheimer solves per occupied state, the static form of \cref{eq:5.4-sth}. The first, $|\psi^{(1)}_v\rangle = -(H_{\mathrm{KS}} - \varepsilon_v)^{-1}\op{P}^{\mathrm{c}}\,\partial_{\hat{\vect{q}}}V|\psi_v\rangle$, is the ordinary phonon response; the second, $|\psi^{(2)}_v\rangle = (H_{\mathrm{KS}} - \varepsilon_v)^{-1}\op{P}^{\mathrm{c}}|\psi^{(1)}_v\rangle$, uses $|\psi^{(1)}_v\rangle$ as its source. Expanding both in conduction states $|c\rangle$ gives one power of $\Delta_{cv}^{-1}$ per solve, so
\begin{equation}
  \mathcal{I}_2 = 2\sum_v\langle\psi^{(1)}_v|\psi^{(1)}_v\rangle , \qquad
  \mathcal{I}_3 = 2\sum_v\langle\psi^{(1)}_v|\psi^{(2)}_v\rangle , \qquad
  \mathcal{I}_4 = 2\sum_v\langle\psi^{(2)}_v|\psi^{(2)}_v\rangle ,
  \label{eq:E.18-sth}
\end{equation}
with the factor 2 for spin and the sums over occupied bands and $\vect{k}$. These are sums over Kohn--Sham pairs. The interacting drag mass at $\mathcal{X}$ = RPA is the $\omega^2$ coefficient of the on-site electronic force constant at imaginary frequency (\cref{sec:12.5}), and a fit in $\omega^2$ through five points between 0.1 and \SI{0.5}{eV} returns it. The difference between the pair sum and the RPA slope is the many-body theory uncertainty of \cref{sec:8.8}. We compute $\delta M_\kappa$ both for the $\vect{Q} = 0$ sublattice displacement and on site (the Brillouin-zone average over an $8^3$ grid of $\vect{Q}$), with the inputs of the shared harmonic base of \cref{app:E.17}. The pseudopotential freezes the core, whose electrons move rigidly with the nucleus; their contribution, $N_{\mathrm{core}}m_e$ to $\delta M_\kappa$, is added exactly, and the core part of $P_{\mathrm{Mig}}$ is added shell by shell from the atomic calculation of \cref{sec:12.2}.

The trace rule tests the result. Displacing the whole cell rigidly drags every electron, so the sum over all sites of the $\vect{Q} = 0$ tensor equals $N_e m_e$ per cell. The trace rule holds only with the core term included, and it tests the $\vect{k}$ grid and the band completeness of \cref{eq:E.18-sth} at once. $\delta M_\kappa$ must change by less than \SI{1}{\percent} from $12^3$ to $16^3$ $\vect{k}$ points. Each perturbation needs one Sternheimer solve beyond those of the shared DFPT run, so the cost is about twice that run: some 200 core-hours for silicon, 30 for germanium and 40 for gallium arsenide, at a few gigabytes of memory, and 2.4 times as much for the $16^3$ test. The Kohn--Sham gap of gallium arsenide, \SI{0.5}{eV} in PBE against \SI{1.52}{eV} measured, puts spurious resonances below the real gap and inflates $\mathcal{I}_3$ and $\mathcal{I}_4$, which carry $\Delta^{-3}$ and $\Delta^{-4}$; a scissor or a hybrid functional is required. Germanium is worse: PBE at the measured lattice constant puts the conduction state at $\mathrm{\Gamma}$ only \SI{0.047}{eV} above the valence maximum, so the inflation of $\mathcal{I}_3$ and $\mathcal{I}_4$ is far larger, and a scissor keeps the wrong valley order. The germanium drag masses, like the electronic piece of the germanium Migdal rate below, are computed on germanium bands that pass the check of \cref{app:E.14}. In a polar crystal the $\vect{Q} = 0$ response must be taken at zero macroscopic field, without the nonanalytic term of \cref{app:E.4}, because the drag mass belongs to the transverse displacement. The results replace the table-based on-site values of \cref{sec:12.5} and fix the sign of the off-site coupling.

\subsection{The crystal Migdal rate}

The Migdal rate of a crystal is set by the Migdal entry $S_{\kappa\kappa}(\vect{q};\omega_l,\omega_e)$ of the joint response, \cref{eq:12.1-Skk}, a density in two frequencies, in \si{eV^{-2}} per cell. The Migdal entry is a convolution of two pieces: the density response $S^{ee}_{\vect{G}\vect{G}'}(\vect{p},\omega_e)$ at the momentum $\vect{p}$ that the boost gives the electrons, with local fields, and the multiphonon or impulse response of the struck nucleus. We want the Migdal entry for silicon and germanium with $p \le \SI{30}{keV}$ and $\omega_e \le \SI{100}{eV}$. On $\mathfrak{T}$ the calculation is at the leading nonadiabatic order with $\mathfrak{M}$ = sudden. It uses the Condon approximation with the lattice in its ground vibrational manifold, at $\mathcal{X}$ = RPA with local fields, and with $\mathcal{K}$ the multiphonon or impulse regime of the lattice tables.

The electronic piece is all-electron: GPAW with PAW reconstruction, or Yambo~\cite{sangalli2019many} and BerkeleyGW~\cite{deslippe2012berkeleygw} with the 2p shell of silicon and the 3d shell of germanium in the valence; the deeper shells enter one at a time from the atomic calculation of \cref{sec:12.2}. The convolution uses the exact impulse bracket of \cref{sec:12.3} wherever $|\vect{q}\cdot\vect{p}|/(M_\kappa\omega_e) > 10^{-2}$, and the soft limit elsewhere. The matrix is not computed on the Cartesian $50^3$ grid of $\vect{p}$ per shell that the convolution integrates over; computing it on that grid would take some $10^{19}$ operations. It is computed at the 72 irreducible reduced momenta of a $12^3$ grid, whose Umklapp copies $\vect{p}_{\mathrm{red}} + \vect{G}$ cover the unreduced momenta, and interpolated in $\vect{p}$ for the convolution. Local fields are converged from $|\vect{G}| \le 4$ to $6\,a_0^{-1}$, which is about 300 to 1000 plane waves in silicon. A $10^3\times10^3$ complex matrix takes \SI{16}{MB} per $(\vect{p},\omega_e)$ and the \SI{0.05}{eV} grid has 2000 frequencies, so a single $\vect{p}$ holds \SI{32}{GB}; the frequencies are processed in blocks, contracted with the lattice response on the fly, and only the convolved $S_{\kappa\kappa}$ is stored. Building the independent-particle matrix costs about $3\times10^{12}$ operations per $(\vect{p},\omega_e)$ on the $12^3$ $\vect{k}$ grid, and the whole calculation about $10^4$ core-hours. The same code, fed the same electronic and lattice tables, reports the rate at each truncation of \cref{fig:12-map}, so that the steps between them are differences of one calculation and not of different inputs.

The third electronic moment, summed over shells and valence, must equal the contact-density sum rule of \cref{sec:12.5}. The $\omega_e^{-3}$ moment must reproduce the drag mass of \cref{app:E.18-drag}, and its trace $N_e m_e$. The computed spectrum must respect the bound on $P_{\mathrm{Mig}}$ of \cref{app:E.18-drag}. No measurement of the crystal Migdal effect exists yet; once the tagged-recoil neutron measurement proposed by Adams et al.~\cite{adams2023measuring} has been made on silicon, the computed rate is compared with it at the same recoil energies. The difference between the result with the full matrix and with its diagonal is reported as the local-field uncertainty.

The Migdal term must not be added to a lattice response computed with the bare mass and unit weight (\cref{sec:12.3}): the lattice part loses the weight $P_{\mathrm{Mig}}$ and moves with the mass $M_\kappa + \delta M_\kappa$, or the same recoil is counted twice. A response diagonal in $\vect{G}$ keeps the local fields on the diagonal of $\varepsilon^{-1}$, which lower the silicon probability by a factor 1.35 against a scalar dielectric function (\cref{lim:12.3a}). It drops the elements $\vect{G}\neq\vect{G}'$, which raise the probability again by a factor of about 1.2 to 1.4 in an empirical-pseudopotential model that is not converged in the reciprocal-lattice sum (\cref{sec:12.3-validity}), and this calculation sizes them. A valence-only electronic piece misses the core-dominated $\omega_e^3$ moment entirely. The results enter \cref{sec:12.3} as the crystal Migdal spectra of silicon and germanium.

\subsection{Coherent zero-phonon Migdal response}

The convolution above treats the lattice factor in its incoherent form, one nucleus at a time and diagonal in $\vect{G}$. A zero-phonon Migdal event is a Bragg process instead: the crystal takes a reciprocal-lattice vector $\vect{G}$ as a whole, with weight $\ee^{-2W_\kappa(\vect{G})}$ and the structure factor of the cell, and the electrons take exactly $\vect{p} = \vect{q} - \vect{G}$ (\cref{sec:12.4}). The amplitudes of different $\vect{G}$ interfere, so the rate is a double sum over $\vect{G}$ and $\vect{G}'$ with the cross-site terms $S_{\kappa\kappa'}$ and the full matrix $S^{ee}_{\vect{q}-\vect{G},\vect{q}-\vect{G}'}$. Above $2q_{\mathrm{BZ}}$ the direction-averaged double sum agrees with the incoherent form. Far below $q_{\mathrm{BZ}}$ the boost identity \cref{eq:8.3-boostward} fixes it at $1/14.9$ of the diagonal form in silicon, with the median $\omega_e$ at the plasmon (\cref{sec:12.4}); a model puts the diagonal form within \SI{15}{\percent} of the coherent one only for $q \le \SI{0.75}{keV}$. Between \SI{0.75}{keV} and $2q_{\mathrm{BZ}}$ lie nearly all of the silicon Migdal rate up to \SI{1}{MeV}, \SI{33}{\percent} of the rate at \SI{2}{MeV} and \SI{7}{\percent} at \SI{3}{MeV} (\cref{tab:12-nuisances}), and in this band of $q$ nothing fixes the double sum.

The calculation evaluates the coherent zero-phonon response of silicon from \SI{0.75}{keV} to $2q_{\mathrm{BZ}}$, at $\omega_l = 0$, with every cross-site term and local fields at $\mathcal{X}$ = RPA. Each $\vect{q}$ is its own reduced momentum plus a reciprocal-lattice vector, so the matrices $S^{ee}_{\vect{G}\vect{G}'}(\vect{p}_{\mathrm{red}},\omega_e)$ of the crystal rate serve unchanged; only the head and wings, which are nonanalytic as $\vect{p}_{\mathrm{red}}\to0$, need reduced momenta added on three lines toward $\mathrm{\Gamma}$, at the cost per momentum of the crystal rate. The ion vertex is the all-electron or PAW one of the crystal rate, and the Bragg weights carry $\ee^{-2W_\kappa(\vect{G})}$ from the shared base. The result is reported as the ratio of the coherent to the diagonal rate per bin of $(q,\omega_e)$, averaged over directions on a 590-point Lebedev grid and kept per direction for the daily modulation. The ratio has three tests: as $q$ falls toward \SI{0.75}{keV} it must approach the boost identity, the factor $1/14.9$ and the plasmon median; above $2q_{\mathrm{BZ}}$ it must return the diagonal form; and at every $q$ it must satisfy the partial sum rule \cref{eq:8.6-fsum-k}. The ratio replaces the unsized entry of \cref{tab:12-nuisances} and the Migdal band of \cref{sec:17.1}.

\subsection{Couplings not proportional to mass}

A scalar that couples to nucleus $\kappa$ with strength $c_\kappa$ is absorbed at momentum $k \to 0$ through the acceleration form \cref{eq:12.6-accel}, with the electronic factor $\langle f|\nabla_{\vect{R}_\kappa}V|i\rangle$. If $c_\kappa/M_\kappa$ is the same for every species, \cref{sec:12.6} reduces the amplitude exactly to the electron density and the rate to the loss function. Write $c_\kappa/M_\kappa = \tilde c + \delta\tilde c_\kappa$. The force operator of a point nucleus is $\nabla_{\vect{R}_I}V = \Vol^{-1}\sum_{\vect{p}}\ii\vect{p}\,Z_\kappa v_{\mathrm{C}}(\vect{p})\,\ee^{-\ii\vect{p}\cdot\vect{R}_I}n_e(\vect{p})$, and the sum over sites of $\delta\tilde c_\kappa\,\ee^{\ii\vect{k}\cdot\vect{R}_I}$ forces $\vect{p} = \vect{k} + \vect{G}$. The part of the rate quadratic in $\delta\tilde c_\kappa$ is then
\begin{equation}
  \begin{gathered}
  \varGamma \propto \frac{1}{\Omega^2\omega^4}\sum_{\vect{G}\vect{G}'}[\vect{k}\cdot(\vect{k}+\vect{G})][\vect{k}\cdot(\vect{k}+\vect{G}')]\,
  \Xi_{\vect{G}}(\vect{k})\Xi^*_{\vect{G}'}(\vect{k})\,S^{ee}_{\vect{G}\vect{G}'}(\vect{k}\to0,\omega) ,
  \\
  \Xi_{\vect{G}}(\vect{k}) = \sum_\kappa \delta\tilde c_\kappa Z_\kappa v_{\mathrm{C}}(\vect{k}+\vect{G})\,\ee^{-\ii\vect{G}\cdot\boldsymbol{\tau}_\kappa} ,
  \end{gathered}
  \label{eq:E.18-X}
\end{equation}
and every Umklapp term is $O(k)$, like the $\vect{G} = 0$ term. The part $\tilde c$ common to all species has the amplitude $-\tilde c\,m_e\langle f|n_e(\vect{k})|i\rangle$ of \cref{eq:12.6-mass}. Interference between this amplitude and the amplitude of $\delta\tilde c_\kappa$ adds to the rate a cross term linear in $\delta\tilde c_\kappa$, with the elements $S^{ee}_{\vect{0}\vect{G}}$ and the weights $[\vect{k}\cdot(\vect{k}+\vect{G})]\,\Xi^*_{\vect{G}}(\vect{k})$. The $\vect{G} = 0$ terms of the two amplitudes add to one coefficient of $n_e(\vect{k})$, and the diagonal element $S^{ee}_{\vect{0}\vect{0}}$ enters only through the loss-function rate \cref{eq:12.6-absrate} with that coefficient. The effect of $\delta\tilde c_\kappa$ beyond the coefficient of $n_e(\vect{k})$ therefore lies in the off-diagonal elements $S^{ee}_{\vect{0}\vect{G}}$ and $S^{ee}_{\vect{G}\vect{G}'}$ at vanishing reduced momentum. For equal $\delta\tilde c$ on both sites of the diamond structure, $\Xi_{\vect{G}}$ is proportional to the basis structure factor, which vanishes for (200) and has $|S^{\mathrm{bas}}|^2 = 2$ and 4 for (111) and (220).

The matrix $S^{ee}_{\vect{G}\vect{G}'}(\vect{k}\to0,\omega)$ is needed for $|\vect{G}| \le 6(2\pi/a)$, about 230 vectors in silicon, and $\omega$ from $\Eg$ to \SI{100}{eV} plus the core edges, contracted with \cref{eq:E.18-X} and with the weights of the cross term for each coupling model. On $\mathfrak{T}$ the calculation is at $\mathcal{X}$ = RPA with local fields (then the adiabatic local-density kernel), at clamped nuclei because $\omega \gg \omega_{\mathrm{max}}$. Yambo or BerkeleyGW give $\chi_{\vect{G}\vect{G}'}(\vect{q}\to0,\omega)$ directly. A diagonal dielectric matrix would miss every Umklapp term, so the local-field set is enlarged until $\varGamma$ and the cross term change by less than \SI{2}{\percent}, on a $12^3$ grid; each frequency needs under \SI{1}{MB}, and the cost is about $10^2$ core-hours. The $\vect{G} = 0$ term uses the full $S^{ee}$, which already contains the factor $1/|\varepsilon|^2$, and screening it again would double count. The calculation is accepted when the (200) channel vanishes to numerical noise for equal $\delta\tilde c$, and when for $c_\kappa \propto M_\kappa$, where $\delta\tilde c_\kappa$ is only the mass defect of order $10^{-3}$, the total rate reduces to the loss-function rate \cref{eq:12.6-absrate}. The result enters \cref{sec:12.6}.

\begin{table}[t]
  \centering
  \caption{Nonadiabatic calculations of silicon (and germanium, gallium arsenide or diamond where named): the computed quantity, the main convergence parameters, the acceptance test and the cost. $N_e$ is the number of electrons per cell; the cost is in core-hours.}
  \label{tab:E18-params}
  \small
  \setlength{\tabcolsep}{4pt}
  \begin{tabular}{p{0.17\textwidth}p{0.29\textwidth}p{0.25\textwidth}p{0.17\textwidth}}
    \toprule
    quantity & grids and parameters & acceptance test & cost \\
    \midrule
    $\delta M_\kappa$, $\mathcal{I}_2$--$\mathcal{I}_4$ & shared base (\SI{60}{Ry}); $\vect{k}$: $12^3\to16^3$ (\SI{1}{\percent}); $8^3$ $\vect{Q}$; $\ii\omega$ = 0.1--\SI{0.5}{eV} & trace $N_e m_e$ per cell & $2\times$ base: 200 (Si), 30 (Ge), 40 (GaAs) \\
    $S_{\kappa\kappa}(\vect{q};\omega_l,\omega_e)$ & $p \le \SI{30}{keV}$ from 72 reduced $\vect{p}$; $\omega_e \le \SI{100}{eV}$ in \SI{0.05}{eV}; $|\vect{G}| \le 4\to6\,a_0^{-1}$; $\vect{k}$: $8^3\to12^3$ & third moment, drag mass, $P_{\mathrm{Mig}}$ bound; neutron data when measured & $10^4$ \\
    coherent zero-phonon ratio & \SI{0.75}{keV} to $2q_{\mathrm{BZ}}$; matrices of the row above plus reduced $\vect{p}$ toward $\mathrm{\Gamma}$; 590 directions & boost limit ($1/14.9$); diagonal form above $2q_{\mathrm{BZ}}$; \cref{eq:8.6-fsum-k} & per added $\vect{p}$ as above \\
    $S^{ee}_{\vect{G}\vect{G}'}(\vect{k}\to0,\omega)$ & $|\vect{G}| \le 6(2\pi/a)$; $\vect{k}$: $12^3$ & (200) zero; mass-proportional limit & $10^2$ \\
    \bottomrule
  \end{tabular}
\end{table}
